%
% Informazioni su:
%   Vision-Language Models,
%   CLIP,
%   BioCLIP,
%   MedCLIP,
%   embeddings,
%   latent space.
%
\chapter{Vision-Language Models}

\section{Introduzione ai Vision-Language Models}

I \textit{Vision-Language Models} (VLM) sono modelli progettati per elaborare e comprendere simultaneamente due differenti modalità di informazione:

\begin{itemize}
    \item la componente visiva, costituita da immagini;
    \item la componente linguistica, costituita da testo naturale.
\end{itemize}

L'obiettivo principale di questi modelli è creare una connessione semantica tra contenuto visivo e descrizioni testuali.

In ambito medico, questo significa associare immagini diagnostiche, come radiografie o risonanze magnetiche, a referti clinici e descrizioni mediche.

Ad alto livello, il funzionamento di un VLM può essere paragonato al comportamento di un medico durante una diagnosi:

\begin{itemize}
    \item osserva immagini cliniche;
    \item legge informazioni testuali;
    \item collega pattern visivi a concetti medici.
\end{itemize}

Un Vision-Language Model cerca quindi di apprendere automaticamente queste associazioni tra elementi visivi e linguaggio clinico.

Ad esempio, dato un esame radiografico toracico, un VLM può determinare se l'immagine risulti semanticamente associata a descrizioni come:

\begin{itemize}
    \item ``versamento pleurico'';
    \item ``polmoni sani'';
    \item ``opacità polmonare''.
\end{itemize}

\section{Architettura Generale dei VLM}

La maggior parte dei moderni Vision-Language Models utilizza un'architettura composta da due rami distinti:

\begin{itemize}
    \item un \textbf{Image Encoder};
    \item un \textbf{Text Encoder}.
\end{itemize}

\subsection{Image Encoder}

L'\textit{Image Encoder} ha il compito di processare immagini e trasformarle in rappresentazioni numeriche compatte (embeddings).

Tra le architetture maggiormente utilizzate troviamo:

\begin{itemize}
    \item Vision Transformers (ViT);
    \item ResNet;
    \item CNN profonde.
\end{itemize}

L'encoder estrae pattern visivi rilevanti dall'immagine e li converte in una rappresentazione latente.

\subsection{Text Encoder}

Il \textit{Text Encoder} elabora invece il testo naturale associato alle immagini.

Tipicamente vengono utilizzate architetture Transformer, tra cui:

\begin{itemize}
    \item BERT;
    \item BioBERT;
    \item ClinicalBERT.
\end{itemize}

Anche il testo viene trasformato in una rappresentazione numerica \textbf{nello stesso} spazio latente dell'immagine.

\section{Contrastive Learning}

Il meccanismo di training più diffuso nei VLM moderni è il cosiddetto \textit{Contrastive Learning}.

Durante l'addestramento, il modello riceve grandi quantità di coppie corrette composte da:

\begin{itemize}
    \item un'immagine;
    \item il relativo testo o referto clinico associato.
\end{itemize}

Il processo di apprendimento ha due obiettivi principali:

\begin{itemize}
    \item massimizzare la similarità tra immagini e testi semanticamente corretti;
    \item minimizzare la similarità rispetto a coppie errate.
\end{itemize}

In altre parole, il modello impara ad avvicinare matematicamente immagini e descrizioni compatibili e ad allontanare quelle non correlate.

Questo approccio consente la costruzione di uno spazio semantico \textbf{condiviso} tra immagini e linguaggio naturale.

\section{CLIP, BioCLIP e MedCLIP}

\subsection{CLIP}

\textit{CLIP} (\textit{Contrastive Language-Image Pretraining}), introdotto da OpenAI, rappresenta uno dei principali modelli pionieristici nel campo dei Vision-Language Models.

Il modello è stato addestrato su enormi quantità di immagini e testo generici provenienti da internet.

CLIP risulta particolarmente efficace nel riconoscimento di oggetti e scene comuni, ma presenta importanti limitazioni nel dominio medico, poiché non è stato addestrato specificamente su dati clinici.

\subsection{MedCLIP e BioCLIP}

Per superare tali limitazioni, sono stati sviluppati modelli specializzati per il dominio biomedicale e medico, tra cui:

\begin{itemize}
    \item MedCLIP;
    \item BioCLIP.
\end{itemize}

Questi modelli mantengono la stessa architettura generale di CLIP, ma vengono addestrati o ri-addestrati utilizzando dataset medici specializzati.

Uno dei dataset più utilizzati è:

\begin{itemize}
    \item \textit{MIMIC-CXR}, contenente centinaia di migliaia di radiografie toraciche associate ai relativi referti clinici.
\end{itemize}

Grazie a questo addestramento specializzato, tali modelli acquisiscono una rappresentazione molto più accurata della semantica clinica.

\section{Embeddings e Spazio Latente}

Uno dei concetti fondamentali nei Vision-Language Models riguarda gli \textit{embeddings} e lo \textit{spazio latente}.

\subsection{Embeddings}

Come ben sappiamo, quando un'immagine o un testo viene processato dall'encoder corrispondente, il risultato non è direttamente interpretabile dall'essere umano, ma viene trasformato in un vettore numerico che prende il nome di \textit{embedding}.

Un embedding può essere rappresentato come:

\[
[0.45, -1.2, 0.88, \dots]
\]

Ogni embedding costituisce una rappresentazione matematica del contenuto semantico dell'immagine o del testo.

\subsection{Spazio Latente}

Gli embeddings generati dal modello vengono collocati all'interno di uno spazio matematico multidimensionale chiamato \textit{latent space}.

In questo spazio:

\begin{itemize}
    \item immagini semanticamente simili risultano vicine tra loro;
    \item testi semanticamente simili risultano vicini tra loro;
    \item Ma soprattutto, immagini e testi associati correttamente condividono regioni simili dello spazio latente.
\end{itemize}

La caratteristica fondamentale dei VLM moderni consiste nella costruzione di uno \textbf{spazio latente condiviso} tra immagini e testo.

Ad esempio:

\begin{itemize}
    \item l'embedding di una radiografia contenente una frattura;
    \item l'embedding della frase ``Frattura scomposta del femore destro'';
\end{itemize}

tenderanno a occupare posizioni molto vicine all'interno dello spazio latente.

\section{Limiti dei Latent Spaces}

Nonostante l'efficacia dei Vision-Language Models, le rappresentazioni latenti risultano estremamente difficili da interpretare.

Le principali problematiche riguardano:

\begin{itemize}
    \item l'elevata dimensionalità delle rappresentazioni;
    \item la natura polisemantica dei neuroni;
    \item la mancanza di interpretabilità umana diretta.
\end{itemize}

Osservando singoli valori all'interno di un embedding, non è possibile comprendere chiaramente quale concetto medico venga rappresentato.

Un determinato neurone potrebbe codificare simultaneamente:

\begin{itemize}
    \item caratteristiche anatomiche;
    \item pattern patologici;
    \item texture visive;
    \item informazioni statistiche apprese durante il training.
\end{itemize}

Di conseguenza, sebbene i modelli producano risultati estremamente accurati, il loro processo decisionale rimane opaco.

Questo rappresenta uno dei principali problemi nell'utilizzo clinico dei sistemi di Intelligenza Artificiale, poiché in medicina è fondamentale comprendere non solo \textit{cosa} il modello predica, ma anche \textit{perché} abbia preso una determinata decisione.